\documentclass[11pt]{article}
\usepackage[margin=1in]{geometry}
\usepackage{amsmath}
\usepackage{array}
\usepackage{booktabs}
\usepackage{longtable}

\title{Memo 01: Mass-Conservation Review of \texttt{conc\_a} and \texttt{conc\_d}}
\author{Codex review for FVCOM-MP floc diagnostics}
\date{\today}

\begin{document}
\maketitle

\section{Purpose}
This memo summarizes the present numerical treatment of the floc-enabled plastic tracers
\texttt{conc\_a} (aggregated microplastic) and \texttt{conc\_d} (disaggregated microplastic)
in \texttt{mod\_plast.F}, with a focus on mass conservation when the interaction term is active.
It also records the leading hypotheses for the observed mass loss.

\section{Present Numerical Treatment}
\subsection{State Variables}
With \texttt{PLAST\_FLOC = T}, the water-column plastic state is split into two tracers:
\[
C = C_a + C_d,
\]
where $C_a$ is stored in \texttt{conc\_a} and $C_d$ is stored in \texttt{conc\_d}. The bed
state is represented by the total areal mass \texttt{mass} together with split bookkeeping
\texttt{mass\_a} and \texttt{mass\_d}.

\subsection{Operator Sequence in the Current Code}
The actual chronology in \texttt{Advance\_plast} is summarized below.

\begin{longtable}{>{\raggedright\arraybackslash}p{0.20\textwidth} >{\raggedright\arraybackslash}p{0.72\textwidth}}
\toprule
Checkpoint & Current treatment \\
\midrule
\endfirsthead
\toprule
Checkpoint & Current treatment \\
\midrule
\endhead
Step start &
The code rebuilds the combined concentration \texttt{conc = conc\_a + conc\_d}, and then enters
the bed-exchange operators. \\
After deposition &
Settling is solved separately for \texttt{conc\_a} and \texttt{conc\_d}. Aggregated plastic
uses a floc/sediment settling velocity, while disaggregated plastic uses the plastic settling
velocity. The stored depositional flux is only the combined scalar
\[
\texttt{dflx} = \texttt{dflx\_a} + \texttt{dflx\_d}.
\]
No separate bed fluxes are retained for the split states. \\
After erosion &
The bed erosion flux \texttt{eflx} is first computed as a total plastic flux. It is then
partitioned back into \texttt{conc\_a} and \texttt{conc\_d} using the instantaneous bottom-layer
ratio
\[
r_b = \frac{\texttt{conc\_a}(k_{bm1})}{\texttt{conc}(k_{bm1})},
\]
with a fallback to purely disaggregated input if the denominator is too small. A similar split
is used for the surface precipitation flux. Both additions are passed through local growth
limiters. \\
After \texttt{Update\_Bottom} &
The total bed mass is updated through the net bed exchange
\[
\Delta M_b \propto \texttt{eflx} - \texttt{dflx},
\]
but the split bed masses are not updated from separate depositional fluxes. Instead, the current
bed split ratio is reused to repartition the updated total mass. In other words, the bed split
evolves by bookkeeping ratio, not by resolved split deposition. \\
After advection/diffusion &
\texttt{conc\_a} and \texttt{conc\_d} are advected independently by two calls to
\texttt{adv\_scal}, then vertically diffused independently by two calls to
\texttt{vdif\_scal}. \\
After kinetics &
The floc interaction step applies the pairwise update
\(C_a^{n+1} = C_a^* + \alpha C_d^* - \beta C_a^*\) and
\(C_d^{n+1} = C_d^* - \alpha C_d^* + \beta C_a^*\),
with $\alpha = \min(\lambda_a \Delta t, 1)$ and $\beta = \min(\lambda_d \Delta t, 1)$.
Algebraically, this step preserves $C_a + C_d$ pointwise before later clipping. \\
After clamps/limiters &
Two non-conservative corrections remain in the floc path: an upper cap
\texttt{cnew\_a > 100} or \texttt{cnew\_d > 100}, and the final nonnegative clip
\texttt{conc\_a = max(cnew\_a, 0)} and \texttt{conc\_d = max(cnew\_d, 0)}. The final combined concentration is then rebuilt
from the clipped split fields. \\
\bottomrule
\end{longtable}

\subsection{Initialization and Restart Notes}
\begin{itemize}
\item Initial water-column concentration is split by \texttt{INI\_CONC\_AG\_RATE}.
\item Initial bottom mass is now split consistently by \texttt{INI\_MASS\_AG\_RATE} for both
constant and static bottom-mass initialization paths.
\item Hot-start support still restores only the combined plastic concentration in the current
restart reader, so floc-state restart consistency remains a concern.
\end{itemize}

\section{Conservation Expectations}
If external sources and open-boundary forcing are disabled, the following operators should be
internally conservative in total plastic mass:
\begin{itemize}
\item separate settling plus bed update, provided the water-column loss exactly matches the bed gain;
\item separate erosion plus bed update, provided the bed loss exactly matches the water-column gain;
\item the pairwise aggregation/disaggregation kinetics, because the gain in one split state is the
loss in the other;
\item interior advection and diffusion, up to the conservation properties of the transport scheme.
\end{itemize}

Therefore, any drift in the internal total
\[
\text{water total} + \text{bed total}
\]
should point to clipping, limiter truncation, split-state bookkeeping inconsistency, or an
operator mismatch between the water-column and bed updates.

\section{Leading Hypotheses for the Observed Mass Loss}
\subsection*{1. Floc-only hard caps and clips are explicit sinks or gains}
The upper cap on \texttt{cnew\_a} and \texttt{cnew\_d} removes positive mass directly. The
nonnegative clip removes negative values, which acts as a positive mass correction.

\subsection*{2. Erosion and surface growth limiters can decouple bed and water updates}
The erosion and surface additions are limiter-controlled in the water column, but the bed update
still uses the full total fluxes \texttt{eflx} and \texttt{dflx}. If the limiter truncates the
water-column increment while the bed still loses the full amount, the difference appears as an
artificial sink.

\subsection*{3. Reuse of one \texttt{cflx} array across both tracer solves may distort floc-mode OBC handling}
Both \texttt{conc\_a} and \texttt{conc\_d} currently call \texttt{adv\_scal} with the same
\texttt{cflx} storage. This may be harmless for strictly interior transport, but it is a plausible
source of distorted boundary-flux bookkeeping when open boundaries are active.

\subsection*{4. Bed split bookkeeping is not driven by split deposition}
The total bed mass is updated with the combined net bed flux, but \texttt{mass\_a} and
\texttt{mass\_d} are then reconstructed from a ratio rather than from resolved split deposition and
erosion fluxes. That makes the split bed inventory diagnostic rather than prognostically exact, and
it can hide where the mismatch is generated.

\section{Suggested Diagnostic Interpretation}
A future diagnostic pass should report domain-integrated values for:
\begin{itemize}
\item water \texttt{a}, water \texttt{d}, water total;
\item bed total, bed \texttt{a}, bed \texttt{d};
\item checkpoint-to-checkpoint totals at:
\texttt{step\_start}, \texttt{after\_deposition}, \texttt{after\_erosion},
\texttt{after\_update\_bottom}, \texttt{after\_advdiff}, \texttt{after\_kinetics}, and
\texttt{after\_clamps};
\item explicit correction counters for upper-cap loss, nonnegative-clip gain, erosion-limiter
truncation, surface-limiter truncation, and bed-split reconciliation.
\end{itemize}

For an internal-isolation test with no point sources and no open-boundary forcing, the most useful
checks are:
\begin{enumerate}
\item the total change from \texttt{after\_advdiff} to \texttt{after\_kinetics}, which should be
near zero if the kinetic exchange is behaving conservatively;
\item the total change from \texttt{step\_start} to \texttt{after\_clamps}, which should be
explained by the reported correction counters.
\end{enumerate}

\section{Bottom Line}
The interaction step itself is algebraically conservative; the stronger candidates for significant
mass loss are the floc-only truncation operators and the mismatch between total bed fluxes and
split-state water-column limiters/bookkeeping. These checkpoints are the first places to inspect in
a later diagnostic pass.

\end{document}
